\documentclass[11pt]{article}
\usepackage[margin=1in]{geometry}
\usepackage{amsmath,amssymb,amsthm,enumitem,hyperref}
\hypersetup{colorlinks=true,linkcolor=blue,urlcolor=blue,citecolor=blue}
\newtheorem{theorem}{Theorem}
\newcommand{\Res}{\operatorname{Res}}
\title{Linearity in Federbush's Genius Conjecture\\\large MathDB 345451}
\author{Principia Math In Progress}
\date{September 18, 2026}
\begin{document}
\maketitle
\begin{abstract}
A formal change of variable absorbs the exponential factor independent
of $y$ in Federbush's coefficient identity. The resulting functions are
coefficients of a composition with the original series and are therefore
linear in its coefficients. An explicit Lagrange coefficient formula
is obtained at the same time.
\end{abstract}

\section*{Task 1. Full solution}
\subsection*{The source's identity}
Fix a positive integer $p$ and set
\[
 U(x)=\sum_{k=1}^p u_kx^k,\quad u_1=1,\qquad
 D(x)=\sum_{k=1}^p d_kx^k,\qquad
 A(x)=e^{D(x)}=\sum_{k\ge0}a_kx^k.
\]
Assume $a_p=0$, as in equation (1) of \cite{fed}. Theorem 1.1 of
that paper defines $F_1=1,F_2,\ldots,F_p$ uniquely by
\begin{equation}\label{eq:target}
 [x^p]A(x)e^{yU(x)}
       =[t^p]\exp\left(y\sum_{i=1}^pF_it^i\right).
\end{equation}
Conjecture 1 asks whether each $F_i$ is linear in the $u_k$.
Here $u_1=1$ is included among the coefficients: equivalently, the
dependence on $u_2,\ldots,u_p$ is affine, with no products of those
variables. We prove the stronger explicit formula below.

\begin{theorem}\label{thm:main}
Define the polynomial
\begin{equation}\label{eq:W}
 W_p(x)=\sum_{k=0}^{p-1}\frac{p}{p-k}a_kx^k.
\end{equation}
Then for every $1\le i\le p$,
\begin{equation}\label{eq:formula}
 \boxed{\quad F_i=\frac1i[x^{i-1}]U'(x)W_p(x)^{i/p}
       =\sum_{k=1}^i\frac{k}{i}u_k
                   [x^{i-k}]W_p(x)^{i/p}.\quad}
\end{equation}
In particular, Conjecture 1 of \cite{fed} holds for every $p$.
\end{theorem}
All fractional powers are formal binomial series with constant term
$1$. The coefficients $a_k$ depend only on the $d_k$, so
\eqref{eq:formula} displays the asserted linear dependence directly.

\subsection*{Absorbing the weight by a change of variable}
Work over the characteristic-zero coefficient ring generated by the
$u_k,d_k$, subject to $a_p=0$. Since $a_0=1$, define
\[
 \phi(x)=xW_p(x)^{-1/p}=x+O(x^2).
\]
It has a unique compositional inverse $\psi(t)=t+O(t^2)$.
Differentiating $\phi(x)^{-p}=x^{-p}W_p(x)$ gives the exact identity
\begin{align}
 \frac{\phi'(x)}{\phi(x)^{p+1}}
 &=x^{-p-1}\left(W_p(x)-\frac{x}{p}W_p'(x)\right)\notag\\
 &=x^{-p-1}\sum_{k=0}^{p-1}a_kx^k.\label{eq:jacobian}
\end{align}
For a formal Laurent series, $\Res_x$ denotes its coefficient of
$x^{-1}$. Let $H(x)$ be any formal power series. Because $a_p=0$,
terms of $A(x)$ beyond those in \eqref{eq:jacobian} have degree at
least $p+1$ and cannot contribute to the residue below. The usual
formal change-of-variable rule consequently gives
\begin{align}
 [x^p]A(x)H(x)
 &=\Res_x H(x)\frac{\phi'(x)}{\phi(x)^{p+1}}\notag\\
 &=\Res_t\frac{H(\psi(t))}{t^{p+1}}
  =[t^p]H(\psi(t)).\label{eq:transport}
\end{align}
This rule can be verified on Laurent monomials: for exponents other
than $-1$ it follows because a derivative has zero residue; for
$t^{-1}$, the residue of $\phi'/\phi$ is $1$. Thus the argument is
purely formal and does not require analytic convergence.

Apply \eqref{eq:transport} with $H(x)=e^{yU(x)}$ and set
\[
 \widetilde F_i=[t^i]U(\psi(t)),\qquad 1\le i\le p.
\]
Then $\widetilde F_1=1$, and the truncation of $U(\psi(t))$ after
degree $p$ cannot affect a coefficient of degree $p$ in its exponential.
Hence the $\widetilde F_i$ satisfy \eqref{eq:target}.

To check uniqueness directly, suppose $F_2,\ldots,F_{i-1}$ are
known. The coefficient of $y^{p-i+1}$ on the right of
\eqref{eq:target} is
\[
 \frac{F_i}{(p-i)!}
       +\text{a polynomial in }F_2,\ldots,F_{i-1}.
\]
Indeed, a contribution involving $F_i$ must consist of that single
degree-$i$ factor and $p-i$ degree-one factors. All other contributions
use only lower indices. This determines $F_i$ successively for
$i=2,\ldots,p$. The leading coefficient in $y$ is already $1/p!$,
and the constant coefficient is $a_p=0$. Therefore
$F_i=\widetilde F_i$. For $p=1$ the assertion is just $F_1=1$.

\subsection*{The explicit coefficient formula}
Formal residue substitution and integration by parts give
\begin{align*}
 [t^i]U(\psi(t))
 &=\Res_x U(x)\frac{\phi'(x)}{\phi(x)^{i+1}}\\
 &=\frac1i\Res_x\frac{U'(x)}{\phi(x)^i}
  =\frac1i[x^{i-1}]U'(x)W_p(x)^{i/p}.
\end{align*}
The integration-by-parts step uses
$\Res_x\frac{d}{dx}(U(x)\phi(x)^{-i})=0$.
This is the classical Lagrange--B\"urmann coefficient identity in
the present coordinates. It proves \eqref{eq:formula} and the
conjecture.\hfill$\square$

As a check, for $p\ge2$ the formula gives
\[
 F_2=u_2+\frac{d_1}{p-1}.
\]
The displayed constant is the coefficient of $u_1=1$; it explains
the affine wording when only $u_2,\ldots,u_p$ are treated as free.
No constraint or hypothesis from the source has been removed.

\subsection*{Attribution and verification}
The setting, unique functions, and conjecture are Federbush's
\cite{fed}; its Bell-polynomial formulation is credited there to
David Williams. Formal residue substitution and Lagrange--B\"urmann
inversion are classical tools. The present application supplies the
coordinate change and the resulting formula for the conjectured
functions.

The verifier uses exact rational arithmetic for three parameter
profiles at each $1\le p\le12$. It checks the residue Jacobian,
the complete identity as a polynomial in $y$, a separate triangular
reconstruction, compositional inversion, linearity with $u_1$ fixed,
and sensitivity to the condition $a_p=0$. The finite checks
supplement the proof for arbitrary $p$.

All six pages of the sole current primary version were read. The
metadata, bounded public searches, and checked claim indexes did not
identify a journal final or a prior full target claim as of September
18, 2026. Review is by the authoring agent only; no independent peer
review or proof-assistant certification is claimed.

\section*{Tasks 2--5. Proposed follow-up directions}
\begin{enumerate}[label=\textbf{Task \arabic*.},start=2,leftmargin=*]
\item Classify when the coefficient transformation in
\eqref{eq:formula} preserves positivity, integrality, or a natural
combinatorial interpretation.
\item Describe composition and inversion laws for successive weight
changes $A(x)$, including the role of the coefficient constraint.
\item Extend the construction to several variables or other
coefficient functionals, and identify the obstructions to an analogous
linear transformation.
\item Study growing-index regimes with $i/p$ bounded away from zero,
where fixed-index expansions need not be uniform.
\end{enumerate}
\begin{thebibliography}{9}
\bibitem{fed} P. Federbush, \emph{The Genius Conjectures (via Bell
Polynomials)}, \href{https://arxiv.org/abs/2002.03814v1}{arXiv:2002.03814v1},
2020. See Sections 1--3, especially Conjecture 1 on page 2.
\end{thebibliography}
\end{document}
